% !TEX root = ../libretto.tex
% Sezione: Canti mariani
\momento{Canti mariani}
\begin{songs}{titoli}

% Fonte: LibrettoAccordi.pdf (Canti Coro), M2.
% Nell'elenco: Canto Mariano, Ordinario.
\beginsong{Ave, Maria}[by={M. Balduzzi, C. Casucci (CD Verbum Panis)}]
\beginchorus
Ave Maria! Ave!
Ave Maria! Ave!
\endchorus
\beginverse
Donna dell'attesa e madre di speranza, ora pro nobis.
Donna del sorriso e madre del silenzio, ora pro nobis.
Donna di frontiera e madre dell'ardore ora pro nobis.
Donna del riposo e madre del sentiero, ora pro nobis.
\endverse
\beginverse
Donna del deserto e madre del respiro, ora pro nobis.
Donna della sera e madre del ricordo, ora pro nobis.
Donna del presente e madre del ritorno, ora pro nobis.
Donna della terra e madre dell'amore, ora pro nobis.
\endverse
\endsong

% Fonte: LibrettoAccordi.pdf (Canti Coro), M12.
% Nell'elenco: Canto Mariano, Ordinario.
\beginsong{Ave Maria}[by={Gen Rosso}]
\beginverse
Ave Maria, piena di grazia,
Il Signore è con te, benedetta fra le donne.
Tu sei la madre di ogni uomo,
tu sei la luce che ci guida nella nostra notte.
Regina del cielo, Madre di Dio,
dolce respiro dell'anima mia.
\endverse
\beginchorus
Ave Maria!
Ave Maria! \rep{2}
\endchorus
\endsong

% Fonte: "11 - Canti Mariani/dell_aurora tu sorgi piu bella ORIGINALE.pdf" (foglio accordi, n. 420).
% Nell'elenco: senza tipologia (Compagnia Aquero).
\beginsong{Dell'aurora tu sorgi più bella}[by={Compagnia Aquero}]
\beginverse
Dell'aurora tu sorgi più bella, coi tuoi raggi a far lieta la terra
e fra gli astri che il cielo rinserra non v'è stella come te.
Gli occhi tuoi son più fondi del mare,
la tua fronte ha il profumo del giglio,
il tuo viso ricorda tuo Figlio, sui tuoi passi nascon fiori.
\endverse
\beginchorus
Bella tu sei qual sole, bianca più della luna
e le stelle più belle non son belle come te. \rep{2}
\endchorus
\beginverse
Ti coronano tutte le stelle, al tuo canto risponderà il vento,
della luna si curva l'argento, si rivolge verso te;
quando tutto dintorno è rovina
e la voce del pianto non tace,
il tuo sguardo riporta la pace, la concordia in fondo ai cuori.
\endverse
\endsong

% Fonte: LibrettoAccordi.pdf (Canti Coro), M6.
% Nell'elenco: Canto Mariano, Ordinario.
\beginsong{Giovane donna}[by={L. Scaglianti, L. Bancolini}]
\beginverse
Giovane donna, attesa dell'umanità,
un desiderio d'amore e pura libertà.
Il Dio lontano è qui vicino a te,
voce e silenzio, annuncio di novità.
\endverse
\beginchorus
Ave Maria! \rep{2}
\endchorus
\beginverse
Dio t'ha prescelta qual Madre piena di bellezza,
ed il suo amore ti avvolgerà con la sua ombra.
Grembo di Dio, venuto sulla terra,
tu sarai madre di un uomo nuovo.
\endverse
\beginverse
Ecco l'ancella che vive della sua parola,
libero il cuore perché l'amore trovi casa.
Ora l'attesa è densa di preghiera
e l'uomo nuovo è qui in mezzo a noi.
\endverse
\endsong

% Fonte: LibrettoAccordi.pdf (Canti Coro), M11.
% Nell'elenco: Canto Mariano, Avvento; l'utente lo vuole tra i Mariani.
\beginsong{L'anima mia magnifica}[by={P. Dargenio, M. Zambuto}]
\beginchorus
L'anima mia magnifica, l'anima mia magnifica.
L'anima mia magnifica, magnifica il Signore. \rep{2}
\endchorus
\beginverse
Perché si è chinato quaggiù e ha guardato alla mia umiltà,
trasformando il mio lamento in una danza.
E in eterno, da ora in poi, canteranno le generazioni
chiameranno il mio nome e per sempre beata sarò.
\endverse
\beginverse
Grandi cose ha compiuto per me il Signore, l'Onnipotente,
che dispiega la Sua pietà su chi lo teme.
La potenza della Sua Maestà ha disperso i superbi di cuore,
rovesciando i potenti e innalzando gli umili.
\endverse
\beginverse
Perché con la Sua carità ha sfamato chi non ha più pane,
rimandando indietro i ricchi a mani vuote.
Ha soccorso il popolo suo, come aveva promesso ad Abramo,
ricordando il Suo Amore fedele per l'umanità.
\endverse
\endsong

% Fonte: LibrettoAccordi.pdf (Canti Coro), M4.
% Nell'elenco: Canto Mariano, Ordinario.
\beginsong{La mia anima canta}[by={Gen Verde}]
\beginchorus
La mia anima canta la grandezza del Signore,
il mio spirito esulta nel mio salvatore.
Nella mia povertà l'Infinito mi ha guardata,
in eterno ogni creatura mi chiamerà beata.
\endchorus
\beginverse
La mia gioia è nel Signore che ha compiuto grandi cose in me,
la mia lode al Dio fedele che ha soccorso il suo popolo
e non ha dimenticato le sue promesse d'amore.
\endverse
\beginverse
Ha disperso i superbi nei pensieri inconfessabili,
ha deposto i potenti, ha risollevato gli umili,
ha saziato gli affamati e aperto ai ricchi le mani.
\endverse
\endsong

% Fonte: LibrettoAccordi.pdf (Canti Coro), M8.
% Nell'elenco: Canto Mariano, Avvento; l'utente lo vuole tra i Mariani.
\beginsong{La mia anima esulta}[by={D. Ricci}]
\beginverse
L'anima mia magnifica il Signore
e il mio spirito esulta in Dio, mio salvatore,
perché ha guardato l'umile sua ancella,
d'ora in poi l'umanità mi chiamerà beata.
Grandi cose, grandi cose ha fatto in me
L'onnipotente, e santo è il nome suo.
La misericordia sua si estenderà
su tutti quelli che lo temono.
\endverse
\beginchorus
L'anima esulta in Dio. La mia anima esulta in Dio.
L'anima esulta in Dio. La mia anima esulta in Dio.
\endchorus
\beginverse
L'anima mia magnifica il Signore
e il mio spirito esulta in Dio, mio salvatore,
perché ha guardato l'umile sua ancella,
d'ora in poi l'umanità mi chiamerà beata.
Ha disteso onnipotente il braccio suo
disperdendo le ambizioni dei superbi, lui.
Giù dai troni ha rovesciato i grandi e
ha posto in alto invece gli umili.
\endverse
\beginverse
L'anima mia magnifica il Signore
e il mio spirito esulta in Dio, mio salvatore,
perché ha guardato l'umile sua ancella,
d'ora in poi l'umanità mi chiamerà beata.
Ha versato beni sopra chi non ha
e i ricchi li ha mandati a mani vuote via.
Ha soccorso noi, suo popolo, perché
lui è fedele nell'amore suo.
\endverse
\endsong

% Fonte: LibrettoAccordi.pdf (Canti Coro), M3.
% Nell'elenco: Canto Mariano, Ordinario.
\beginsong{Madre dolcissima}[by={Gen Rosso}]
\beginverse
Se io fossi poeta, scriverei di te le cose che nessuno ha detto mai.
Dell'arte avessi il dono, io rapirei la luce per poter disegnare gli occhi tuoi.
Se grandi avessi le ali, le scioglierei nel volo per arrivare in alto fino a te.
\endverse
\beginchorus
Madre Dolcissima,
tesoro dei più poveri, coraggio di chi ancora spera. Tu,
Madre Dolcissima, consola queste lacrime, ascolta chi parole più non ha.
\endchorus
\beginverse
Riposo dei più deboli, silenzio di chi dà, fontana per la nostra siccità,
splendore della sera, stella del mattino, neve bianca sopra le città,
regina della pace, vetta dell'amore, radice della nostra vita, tu
\endverse
\beginverse
Dei naufraghi tu l'ancora, del canto poesia, dei nostri sogni unica realtà,
riflesso d'infinito, chiarore delle stelle, immagine del cosmo che sarà,
custode della terra, sorgente fra le rocce, segreto schiuso da ogni fiore, tu,
\endverse
\endsong

% Fonte: LibrettoAccordi.pdf (Canti Coro), M5.
% Nell'elenco: Canto Mariano, Ordinario.
\beginsong{Maria}[by={RnS}]
\beginverse
Ora che sei qui, la tua dolcezza sento dentro me.
Come l'aurora sorgi, risplendi fra noi.
Ora che sei qui, tu figlia e madre, fra le braccia tue
diffondi la tua luce celeste di pace e di amor.
\endverse
\beginchorus
Tu Maria, tu sei dimora per Gesù figlio tuo.
Regina che dal cielo speranza e amore ci dai.
\endchorus
\beginverse
Ora che sei qui, tu dolce madre sono figlio tuo
trasforma il mio pianto nel canto d'amore per te.
\endverse
\beginverse
Ora che sei qui, insegnami ad amare come sai,
radiosa tra le stelle, preghiera ora sei.
\endverse
\endsong

% Fonte: LibrettoAccordi.pdf (Canti Coro), M10.
% Nell'elenco: Canto Mariano, Ordinario.
\beginsong{Mi alma canta}[by={Gen Verde}]
\beginchorus
Mi alma canta, canta la grandeza del Señor,
mi espíritu se estremece de gozo en Dios,
mi Salvador.
\endchorus
\beginverse
Porque miró con bondad
la pequeñez de su servidora.
Porque miró con bondad
la pequeñez de su servidora.
En adelante todas las gentes
me llamarán feliz!
Me llamarán feliz!
Me llamarán feliz!
\endverse
\beginverse
Derribó del trono a los poderosos,
y elevó a los humildes,
colmó de bienes a los hambrientos
y despidió a los ricos con las manos vacías.
Mi alma canta la grandeza del Señor
y mi espíritu se estremece
de gozo en Dios, mi Salvador.
\endverse
\endsong

% Fonte: LibrettoAccordi.pdf (Canti Coro), M9.
% Nell'elenco: Canto Mariano, Ordinario. CD: Giorno senza tramonto.
\beginsong{Regina della benevolenza}[by={M. Balduzzi (CD Giorno senza tramonto)}]
\beginverse
Stella del mare, signora del cielo,
porto sicuro dell'anima nelle tempeste.
Casa dei popoli, donna d'amore,
dolce regina che dona il re dell'universo.
\endverse
\beginchorus
Madre di Dio, prega per noi,
prega, prega, prega per noi.
\endchorus
\beginverse
Mano che veglia sui cuori nella notte buia,
sguardo di benevolenza sopra la città,
aurora chiara nel giorno nuovo che ci aspetta,
mistica rosa fiorita sulle ingiurie della terra.
\endverse
\endsong

% Fonte: LibrettoAccordi.pdf (Canti Coro), M1. C'e' anche "11 - Canti Mariani/RIMANI CON ME.pdf".
% Nell'elenco: Canto Mariano, Ordinario.
\beginsong{Rimani con me}[by={Gen Rosso}]
\beginverse
Maria, tu sei la vita per me,
sei la speranza, la gioia, l'amore: tutto sei.
Maria tu sai quello che vuoi,
sai con che forza d'amore in cielo mi porterai.
\endverse
\beginchorus
Maria, ti do il mio cuore per sempre se vuoi,
tu dammi l'amore che non passa mai.
Rimani con me e andiamo nel mondo insieme.
La tua presenza sarà goccia di paradiso per l'umanità.
\endchorus
\beginverse
Maria, con te sempre vivrò,
in ogni momento giocando, cantando, ti amerò.
Seguendo i tuoi passi in te io avrò
la luce che illumina i giorni e le notti dell'anima.
\endverse
\endsong

% Fonte: LibrettoAccordi.pdf (Canti Coro), M7.
% Nell'elenco: Canto Mariano, Ordinario.
\beginsong{Santa Maria del cammino}[by={R. D'Andrea, J. A. Espinosa}]
\beginverse
Mentre trascorre la vita
solo tu non sei mai;
Santa Maria del cammino
sempre sarà con te.
\endverse
\beginchorus
Vieni, o Madre, in mezzo a noi, vieni Maria quaggiù.
Cammineremo insieme a te
verso la libertà.
\endchorus
\beginverse
Quando qualcuno ti dice:
«Nulla mai cambierà»,
lotta per un mondo nuovo,
lotta per la verità!
\endverse
\beginverse
Lungo la strada la gente
chiusa in se stessa va;
offri per primo la mano
a chi è vicino a te.
\endverse
\beginverse
Quando ti senti ormai stanco
e sembra inutile andar,
tu vai tracciando un cammino:
un altro ti seguirà.
\endverse
\endsong

\end{songs}
